%------------------------------------------------------------------------------%
\begin{question}
  Determine $k$ para que o paralelepípedo determinado por
  \[
  \vec{v}=(1,2,1)
  \]
  \[
  \vec{w}=(2,-1,3)
  \]
  \[
  \vec{z}=(k,1,2)
  \]
  tenha volume igual a $6$
\end{question}

%------------------------------------------------------------------------------%
\begin{resolution}{Exemplo 5 — Resolução}
O volume é
\[
V=
\left|
\begin{vmatrix}
1&2&1\\
2&-1&3\\
k&1&2
\end{vmatrix}
\right|
\]
Expandindo pela primeira linha
\[
V=
\left|
\begin{vmatrix}
-1&3\\
1&2
\end{vmatrix}
-
2
\begin{vmatrix}
2&3\\
k&2
\end{vmatrix}
+
\begin{vmatrix}
2&-1\\
k&1
\end{vmatrix}
\right|
\]
\[
V=
|-5-2(4-3k)+(2+k)|
\]
\[
V=
|6k-11|
\]
Como o volume deve ser $6$
\[
|6k-11|=6
\]
\end{resolution}

%------------------------------------------------------------------------------%
\begin{resolution}{Exemplo 5 — Resolução}
Temos duas possibilidades
\[
6k-11=6
\]
ou
\[
6k-11=-6
\]
Da primeira
\[
6k=17
\]
\[
k=\frac{17}{6}
\]
Da segunda
\[
6k=5
\]
\[
k=\frac56
\]
Portanto os valores possíveis são
\[
k=\frac56
\]
e
\[
k=\frac{17}{6}
\]
\end{resolution}
%------------------------------------------------------------------------------%
